\section{Анализ задачи и обзор существующих решений}

Разработка современного фитнес-приложения, претендующего на роль интеллектуального помощника, требует комплексного анализа двух основных областей: актуальных технологических стеков для мобильной разработки и научно обоснованных принципов построения тренировочного процесса и восстановления. В рамках данного исследования был проведён детальный обзор как технических публикаций и платформенной документации, так и современных методических и научных работ в области физиологии спорта и нейронаук. \cite{refa,refb}

\subsection{Технологический стек и архитектурные подходы}

Основу технической реализации проекта составило изучение официальной документации и рекомендаций для разработчиков на платформе Android. \cite{refc} Современная экосистема Android сместилась в сторону декларативных UI-фреймворков, что подтверждается глубоким анализом возможностей Jetpack Compose. \cite{refd} Этот инструмент не просто упрощает создание интерфейсов, но и способствует построению более отзывчивых и легких в поддержке приложений за счёт реактивной модели данных и композиции компонентов. Изучение архитектурных шаблонов, сопутствующих Compose (таких как Model-View-ViewModel или MVI), выявило их преимущества в управлении состоянием UI и отделении бизнес-логики от кода отображения, что критически важно для долгосрочной поддержки проекта. \cite{refe}

Для организации работы с данными были рассмотрены стандартные решения. Библиотека Room, как часть Android Jetpack, подтвердила свою эффективность в качестве абстракции над SQLite, предоставляя безопасность типов на уровне компиляции, удобные аннотации и поддержку наблюдения за изменениями данных через интеграцию с LiveData или Flow. \cite{reff} Сетевой слой требовал надёжного решения для REST API, и выбор пал на Retrofit, который де-факто является отраслевым стандартом благодаря генерации реализации интерфейсов, встроенной поддержке различных конвертеров (включая Gson и Moshi) и высокой степени кастомизации через интерцепторы. \cite{refg} Особое внимание было уделено изучению паттернов объединения локального кэша (Room) и сетевых данных, включая стратегии синхронизации и обработки офлайн-сценариев. \cite{refh}

Также в обзор попали технологии для асинхронного программирования, где Kotlin Coroutines в сочетании с Flow продемонстрировали значительные преимущества перед традиционными callback-подходами и RxJava в контексте Kotlin-first разработки, предлагая более простой и читаемый код для выполнения фоновых операций. \cite{refi}

\subsection{Методологические основы тренировок и восстановления}

Функциональное наполнение приложения проектировалось с опорой на современные спортивные методологии. Ключевым стал анализ принципа прогрессивной перегрузки, согласно которому для роста физических показателей необходимо систематически увеличивать нагрузку на организм. \cite{refj} Однако исследования подчёркивают, что этот рост нелинеен и зависит от индивидуальных адаптационных возможностей и качества восстановления. Это обосновало отказ от жёстких, универсальных тренировочных планов в пользу системы, которая адаптирует рекомендации по тренировочной нагрузке на основе истории выполнения пользователем упражнений.

Анализ научных работ, посвящённых влиянию сна на физическую и когнитивную производительность, показал, что недостаток или низкое качество сна являются критическими факторами, ограничивающими прогресс. \cite{refk} Они приводят к снижению эффективности восстановления, ухудшению моторного контроля и повышению риска травм. Эти данные сформировали методологическую основу для интеграции ручного трекера восстановления, где сон рассматривается как активный компонент тренировочного процесса, равный по значимости нагрузке.

Таким образом, методология приложения основывается на двух ключевых переменных: тренировочной нагрузке (выполненный объём упражнений) и продолжительности сна. Алгоритм анализирует эти данные по простому принципу, формируя основу для адаптивного планирования. \cite{refj,refk}

\subsection{Анализ существующих решений и рыночных тенденций}

Для определения уникального ценностного предложения был проведён сравнительный анализ популярных фитнес-приложений. Изучались специализированные программы для силового тренинга (Strong, Hevy, Jefit), комплексные фитнес-трекеры (MyFitnessPal, Adidas Training) и встроенные решения от производителей умных часов (Samsung Health, Google Fit). Выявились общие тенденции, определяющие текущие рыночные пробелы.

Большинство приложений сосредоточено на одной из двух функций: либо это цифровой дневник для регистрации подходов и повторений с обширной, но статичной библиотекой упражнений, либо это агрегатор данных из различных источников (шаги, пульс, сон), слабо связанных с конкретными тренировочными целями. Планы тренировок часто представляют собой жесткие, неизменяемые шаблоны, рассчитанные на «среднего» пользователя. В них отсутствует адаптивность, то есть способность автоматически подстраивать сложность и объём нагрузки в ответ на фактические результаты и состояние пользователя. \cite{refl}

Ключевой недостаток — это разрыв между данными о нагрузке (тренировки) и восстановлении (сон, утомление). Даже при наличии трекеров сна эта информация редко используется для модификации последующих тренировочных заданий.

Таким образом, рынок заполнен приложениями, которые действуют либо как пассивный журнал, либо как сборщик разрозненных метрик, но не как активный, интеллектуальный помощник, принимающий решения. \cite{refl} Однако научные обзоры и маркетинговые исследования указывают на растущий спрос со стороны пользователей на более персонализированные и «умные» подходы в сфере здоровья и фитнеса. Пользователи ожидают, что приложение не просто хранит данные, но анализирует их и даёт понятные, обоснованные рекомендации, которые эволюционируют вместе с их прогрессом и текущим состоянием. Это создаёт запрос на системы, способные на простой, но эффективный анализ введённой информации для персонального руководства, что и является основным направлением для разрабатываемого решения.

\subsection{Анализ и выбор средств реализации}

Для достижения поставленных целей и решения задач будут использованы следующие методы анализа и исследования: анализ существующих приложений для отслеживания расписания и API, проектирование пользовательского интерфейса с использованием методов UX/UI дизайна, а также разработка и тестирование программного кода приложения. \cite{refa,refb}

\subsubsection{Обзор существующих языков программирования}

Выбор языка программирования является фундаментальным решением, определяющим архитектуру, производительность и поддерживаемость всего проекта. В экосистеме Android традиционно доминировала Java, однако с 2017 года Kotlin, представленный JetBrains, получил официальную поддержку Google и стал приоритетным языком для платформы. Сравнительный анализ выявил существенные различия в философии и прагматике разработки на этих языках. \cite{refb}

Java, будучи зрелым, объектно-ориентированным языком с огромной экосистемой и проверенной временем JVM, предоставляет стабильность и предсказуемость. Однако для Android-разработки его синтаксис зачастую избыточен. Создание даже простых data-классов требует написания геттеров, сеттеров, методов `equals()`, `hashCode()` и `toString()` вручную или с помощью генераторов IDE, что увеличивает объем кода и потенциальные источники ошибок. Обработка null-значений, частая причина `NullPointerException`, полностью ложится на разработчика, требуя постоянных проверок. Асинхронные операции и работа с потоками, критически важные для отзывчивого UI, реализуются с помощью сложных механизмов вроде `AsyncTask`, `ThreadPoolExecutor` или сторонних библиотек, таких как RxJava, что усложняет код и его отладку.

В противовес этому, Kotlin с самого проектирования задумывался как прагматичный и лаконичный язык, устраняющий boilerplate-код. Data-классы в Kotlin объявляются одной строкой, автоматически получая все стандартные методы. Система null-безопасности, встроенная в систему типов, устраняет целый класс ошибок на этапе компиляции, различая nullable и non-nullable типы. Для асинхронного программирования Kotlin предлагает корутины — легковесные потоки, позволяющие писать асинхронный код в последовательном, императивном стиле, что радикально упрощает его понимание и поддержку. Дополнительные функции, такие как расширения, позволяющие "добавлять" методы к существующим классам, и параметры по умолчанию в функциях, делают код более выразительным и модульным. \cite{refc}

С точки зрения производительности оба языка компилируются в байт-код JVM и демонстрируют сопоставимые результаты в чистом выполнении. Однако производительность разработки и надежность итогового кода в Kotlin выше благодаря упомянутым возможностям. Компилятор Kotlin активно выполняет оптимизации и выявляет потенциальные проблемы (например, неиспользуемые переменные) на ранних этапах. Полная интероперабельность с Java позволяет Kotlin-проекту использовать любые существующие Java-библиотеки, обеспечивая плавную миграцию и доступ к огромной экосистеме.

Сообщество и поддержка стремительно смещаются в сторону Kotlin. Официальная документация Android, новые библиотеки Jetpack (такие как Compose) в первую очередь ориентированы на Kotlin. Инструменты разработки, включая Android Studio, имеют превосходную поддержку Kotlin с продвинутыми рефакторингами и анализом кода. В итоге, несмотря на то, что Java остается надежным выбором для поддержки legacy-проектов, для новой разработки Kotlin предлагает более современный, безопасный, выразительный и эффективный инструмент, напрямую влияющий на скорость разработки и качество кода. \cite{refc,refd}

\subsubsection{Обзор существующих фреймворков для интерфейса}

Эволюция подходов к созданию пользовательского интерфейса в Android прошла путь от жесткой императивной модели с виджетами на основе `View` и `ViewGroup` до современной декларативной парадигмы. Традиционный подход, использовавшийся более десяти лет, основывается на XML-разметке для описания иерархии элементов и отдельного Java/Kotlin кода для управления их состоянием и поведением (логика в `Activity`/`Fragment`). Этот метод, хотя и знаком миллионам разработчиков, имеет присущие ему недостатки: сильная связанность логики и представления, сложность синхронизации состояния UI с данными, необходимость ручного поиска элементов через `findViewById` (позже — View Binding) и громоздкие механизмы обновления сложных интерфейсов.

Jetpack Compose, анонсированный как современный набор инструментов для верстки пользовательского интерфейса, кардинально меняет эту парадигму. Вместо XML, Compose позволяет описывать интерфейс декларативно с помощью функций на Kotlin, которые называются composable-функциями. Это означает, что разработчик описывает, как должен выглядеть UI при определенном состоянии, а система Compose берет на себя ответственность за его отрисовку и обновление при изменении этого состояния. Такой подход устраняет необходимость в ручном управлении виджетами и их свойствах, снижая вероятность ошибок, таких как утечки памяти из-за неправильной привязки к жизненному циклу. \cite{refd}

Ключевое преимущество Compose в производительности — это механизм "интеллектуальной рекомпозиции". Когда изменяется состояние, от которого зависит composable-функция, Compose пересчитывает (рекомпонирует) только те части UI, которые действительно изменились, а не всю иерархию целиком, как это часто бывало в View-системе. Это достигается за счет отслеживания чтения состояния внутри функций. В сочетании с ленивыми компонентами для эффективной работы со списками это приводит к плавной прокрутке и отзывчивости даже на устройствах с ограниченными ресурсами. \cite{refd}

Гибкость и выразительность Compose превосходят возможности XML. Создание сложных, кастомных компонентов с анимациями становится значительно проще. Анимации являются первоклассными компонентами в Compose и могут быть добавлены буквально в одну строку кода. Кроме того, Compose тесно интегрирован с Material Design 3, предоставляя готовые, современные и адаптивные компоненты, которые автоматически следуют рекомендациям по дизайну (цветовая схема, типографика, формы). При этом система остается открытой для создания полностью кастомного дизайна. \cite{refd}

Хотя сообщество и количество обучающих материалов для традиционного подхода по-прежнему огромны, Compose быстро набирает популярность. Google позиционирует его как будущее UI-разработки под Android, активно развивая и добавляя новые возможности. Важно отметить, что Compose может сосуществовать с существующими View-компонентами в одном проекте (через `AndroidView`), что позволяет постепенно мигрировать крупные приложения.

Для тестирования UI Compose предоставляет специальный набор инструментов (`compose-ui-test`), который позволяет писать семантические тесты, взаимодействуя с интерфейсом на высоком уровне (например, "найти элемент с текстом 'Сохранить' и нажать на него"), что делает тесты более стабильными и читаемыми по сравнению с инструментами вроде Espresso для View-системы. \cite{refe}

\subsubsection{Обзор существующих фреймворков для работы с базой данных}

Локальное хранение структурированных данных является критически важным компонентом для любого офлайн-функционального или эффективно кэширующего данные приложения. Нативным решением для Android на протяжении всей истории платформы является SQLite — легковесная, встраиваемая реляционная СУБД. Работа с низкоуровневым SQLite через `SQLiteOpenHelper` и `Cursor` предоставляет разработчику максимальный контроль: можно писать любые, даже самые экзотические SQL-запросы и вручную оптимизировать схемы баз данных. Однако этот подход сопряжен с серьезными издержками. Разработчик вынужден вручную писать большое количество шаблонного кода для создания и обновления таблиц, преобразования объектов приложения в `ContentValues` и обратно, а также для безопасного выполнения транзакций. Любая ошибка в SQL-запросе (опечатка в имени столбца, неверный тип) обнаруживается только во время выполнения приложения (runtime), что может привести к падению. Процесс миграции схемы базы данных при выходе новых версий приложения также требует крайне аккуратной ручной реализации. \cite{reff}

Room Persistence Library, часть Android Jetpack, была создана как абстракция над SQLite, решающая именно эти проблемы. Она реализует популярный паттерн ORM (Object-Relational Mapping), позволяя работать с объектами Kotlin/Java напрямую. Компоненты Room — это аннотированные классы: `@Entity` для определения таблиц, `@Dao` (Data Access Object) для интерфейсов или абстрактных классов, содержащих методы доступа к данным, и `@Database` для описания самой базы. Ключевое преимущество Room — проверка SQL-запросов на этапе компиляции. Компилятор анализирует аннотации `@Query` и сверяет их с объявленными сущностями (`@Entity`). Если запрос ссылается на несуществующее поле или содержит синтаксическую ошибку, компиляция не пройдет, предотвращая множество критических ошибок в рантайме.

С точки зрения чистой скорости выполнения запросов Room и нативный SQLite, как правило, эквивалентны, поскольку Room является тонкой оберткой и в конечном счете генерирует оптимизированный SQL-код. Однако производительность разработки и надежности при использовании Room значительно выше. Библиотека автоматически генерирует код для `@Dao` и работы с `@Database`, избавляя от рутинного кодирования. Встроенная поддержка LiveData и Flow в методах `@Dao` позволяет легко создавать реактивные потоки данных: при изменении информации в таблице все активные наблюдатели (например, UI-компоненты) автоматически получают обновленные данные, что идеально сочетается с архитектурой, основанной на наблюдении за состоянием. \cite{reff}

Гибкость Room, хотя и несколько ограниченная по сравнению с прямым SQL, достаточна для подавляющего большинства сценариев. Она поддерживает сложные запросы с джойнами, вложенными объектами (через `@Relation` и `@Embedded`), кастомные преобразователи типов (`@TypeConverter`). Миграция схемы осуществляется через предоставление явных объектов `Migration`, что делает процесс контролируемым и тестируемым. Для тестирования Room предлагает режим базы данных в памяти (`Room.inMemoryDatabaseBuilder`), что позволяет писать быстрые и изолированные unit-тесты без необходимости работы с файловой системой. \cite{reff}

Комьюнити и поддержка Room исключительно сильны, так как это официальная рекомендация Google для работы с локальной базой данных. Библиотека активно развивается, интегрируется с другими компонентами Jetpack (например, Paging 3) и имеет обширную документацию.

\subsubsection{Обзор существующих фреймворков для работы с сетью}

Сетевой слой — это мост между клиентским приложением и серверной логикой, и его надежность, эффективность и простота использования напрямую влияют на пользовательский опыт. При анализе инструментов для Android фокус сместился с рассмотрения низкоуровневых `HttpURLConnection` на библиотеки более высокого уровня. Ведущими кандидатами в этой области являются OkHttp и Retrofit, причем важно понимать их концепцию: Retrofit не заменяет OkHttp, а строится поверх него, добавляя уровень абстракции. \cite{refg}

OkHttp, разработанный компанией Square, — это мощный HTTP-клиент, реализующий современные стандарты сети. Он берет на себя множество сложных задач: управление пулом соединений (connection pooling) для уменьшения задержки при повторных запросах, прозрачное сжатие GZIP для экономии трафика, кэширование ответов в соответствии с HTTP-спецификацией для работы в офлайн-режиме или ускорения повторных запросов, а также поддержку HTTP/3 и WebSockets. OkHttp предоставляет гибкий механизм интерцепторов (Interceptors), который позволяет встраиваться в цепочку обработки запроса и ответа. Это используется для добавления общих заголовков (например, авторизации), логирования, обработки ошибок или даже модификации запросов на лету. Однако использование "чистого" OkHttp для типичных REST API все еще требует написания значительного объема кода для сериализации объектов приложения в JSON (и обратно) и организации асинхронных вызовов.

Retrofit, также продукт Square, решает эти проблемы, предоставляя декларативный подход к определению сетевых операций. Разработчик описывает API в виде интерфейса на Kotlin, аннотируя методы и их параметры (`@GET`, `@POST`, `@Path`, `@Query`, `@Body` и др.). Retrofit во время выполнения генерирует реализацию этого интерфейса, которая преобразует вызовы метода в HTTP-запросы. Он интегрируется с популярными конвертерами (Gson, Moshi, Jackson) для автоматической сериализации и десериализации данных, что превращает трудоемкий процесс работы с `JSONObject` или `JsonReader` в простую работу с типами данных Kotlin (data-классами). Асинхронность легко достигается за счет возвращаемого типа метода: изначально это был `Call<T>`, теперь же стандартом стала поддержка Kotlin coroutines через `suspend`-функции, что делает сетевой код линейным и легко читаемым. \cite{reff}

Таким образом, связка Retrofit + OkHttp является оптимальной. Retrofit выступает как фасад для удобного описания API, а OkHttp остается своеобразным «двигателем», отвечающим за транспортные операции, безопасность и низкоуровневые оптимизации. Разработчик может сконфигурировать кастомный экземпляр `OkHttpClient` с нужными таймаутами, интерцепторами (например, для добавления логов или индикатора загрузки) и кэшем, а затем передать его в билдер Retrofit. \cite{reff}

Гибкость этой связки крайне высока. Через кастомные `CallAdapter.Factory` можно адаптировать Retrofit для работы с любыми другими механизмами (например, RxJava `Observable`, хотя сейчас это менее актуально). `Converter.Factory` позволяет работать с нестандартными форматами данных (не JSON). Обработка ошибок унифицирована: можно использовать интерцепторы OkHttp для обработки кодов ответа на уровне всего приложения или обрабатывать исключения (например, `HttpException` или сетевые ошибки) непосредственно в слое репозитория при работе с корутинами. \cite{reff}

Тестирование сетевого слоя, построенного на Retrofit, также упрощено. Для модульных тестов можно легко создать mock-реализацию интерфейса API, возвращающую заданные данные. Для интеграционных тестов можно поднять локальный mock-сервер (например, с помощью MockWebServer из состава OkHttp) и проверить, что приложение корректно формирует запросы и обрабатывает ответы. \cite{reff}

\subsubsection{Асинхронное программирование}

В современных мобильных приложениях непрерывно выполняются сетевые запросы, операции с базой данных, обработка больших объёмов данных и обновление пользовательского интерфейса. Если эти задачи будут выполняться в главном потоке (UI-потоке), приложение перестанет реагировать на касания, что приведёт к «зависаниям» и негативному пользовательскому опыту. Асинхронное программирование решает эту проблему, позволяя запускать длительные операции в фоновых потоках без блокировки интерфейса. Как подчёркивает Ёранссон в своей работе, эффективное управление потоками и асинхронными задачами критически важно для мобильных приложений, так как минимизирует затраты ресурсов и сохраняет отзывчивость \cite{refm}.

В экосистеме Android сложилось два основных подхода к организации асинхронности: RxJava и Kotlin Coroutines.

RxJava — библиотека, реализующая принципы реактивного программирования. Она оперирует потоками данных, представленными в виде Observable (или Flowable), и предоставляет богатый набор операторов для трансформации, фильтрации, объединения и агрегации этих потоков. Управление потоками выполнения осуществляется с помощью Schedulers (например, \texttt{Schedulers.io()} для сетевых операций, \texttt{AndroidSchedulers.mainThread()} для UI). RxJava даёт разработчику мощный инструментарий для работы со сложными асинхронными сценариями, включая обработку ошибок, повторные попытки и управление временем. Однако у RxJava есть недостатки: крутая кривая обучения, большое количество абстракций и шаблонного кода, а также повышенное потребление памяти при создании цепочек операторов.

Более современным и предпочтительным для Kotlin-проектов является использование корутин (Coroutines). Корутины — это легковесные потоки, которые не привязаны к операционной системе и управляются самим языком. В отличие от традиционных потоков, создание корутины не требует выделения большого объёма памяти (типичная корутина занимает порядка нескольких килобайт), поэтому их можно создавать тысячами без риска переполнения памяти.

Ключевая концепция корутин — \texttt{suspend}-функции. Такая функция может приостановить своё выполнение в определённой точке, не блокируя поток, в котором она работает, и возобновиться позже, когда результат будет готов. Например, при вызове сетевого запроса корутина «засыпает», освобождая поток для других задач, а когда ответ сервера приходит — «просыпается» и продолжает выполнение с того же места. Это позволяет писать асинхронный код в привычном императивном стиле (последовательные инструкции), без вложенных колбэков или цепочек операторов.

Особый интерес представляет внутреннее устройство корутин. При компиляции каждая \texttt{suspend}-функция преобразуется компилятором Kotlin в конечный автомат (state machine). Каждая точка приостановки (вызов другой \texttt{suspend}-функции) становится состоянием этого автомата. Локальные переменные и состояние сохраняются в специальном объекте, который передаётся между состояниями. Такой подход позволяет эффективно управлять приостановками и возобновлениями без накладных расходов на создание отдельных потоков.
Ключевую роль в этом механизме играет интерфейс \texttt{Continuation}. Каждая \texttt{suspend}-функция неявно получает дополнительный параметр типа \texttt{Continuation<T>}, который передаётся компилятором. Этот интерфейс представляет собой колбэк: он содержит метод \texttt{resume(value: T)} для успешного возобновления и \texttt{resumeWithException(exception: Throwable)} для передачи ошибки. Состояние конечного автомата (текущая метка точки приостановки, локальные переменные) сохраняется в объекте, реализующем \texttt{Continuation}. При вызове \texttt{suspend}-функции текущее продолжение (continuation) передаётся в неё; когда асинхронная операция завершается, вызывается \texttt{resume} или \texttt{resumeWithException}, и выполнение корутины продолжается с того места, где она была приостановлена. Благодаря этому подходу корутины не требуют выделения отдельного потока для каждой приостановки, а вся логика управления состояниями эффективно кодируется в сгенерированной байт-коде.


Кроме того, корутины поддерживают структурированную конкурентность (structured concurrency). Это означает, что все корутины, запущенные в определённой области (например, в рамках жизненного цикла ViewModel), автоматически отменяются при выходе из этой области. Это предотвращает утечки памяти и избавляет от необходимости вручную отслеживать завершение фоновых задач. Для управления потоками выполнения используются диспетчеры (\texttt{Dispatchers}): \texttt{Dispatchers.IO} для операций ввода-вывода (сеть, диск), \texttt{Dispatchers.Default} для CPU-интенсивных задач и \texttt{Dispatchers.Main} для UI.

Корутины тесно интегрированы с другими компонентами Android и Jetpack: \texttt{ViewModel} имеет встроенную поддержку корутин через \texttt{viewModelScope}, Room предоставляет \texttt{suspend}-функции для DAO, а Retrofit поддерживает \texttt{suspend}-функции в своих сервисах. Это делает корутины естественным выбором для современных приложений. Как отмечает Роман Елизаров, корутины позволяют легко интегрировать асинхронные операции в обычный поток выполнения, упрощая код и делая его более читабельным без лишних абстракций \cite{refn}.

Таким образом, для реализации асинхронных операций в разрабатываемом приложении выбраны Kotlin Coroutines в связке с \texttt{Flow} (асинхронные потоки данных, также построенные на корутинах). Это обеспечивает высокую производительность, простоту написания и сопровождения кода, а также бесшовную интеграцию с используемыми библиотеками (Room, Retrofit, Jetpack Compose).

\subsubsection{Выбор средств реализации}

В качестве языка программирования был выбран Kotlin, потому что это современный, лаконичный и мощный язык, который устраняет многие недостатки Java. Kotlin предлагает такие функции, как null-безопасность, корутины, data-классы и расширения, которые значительно упрощают разработку и делают код более читаемым и поддерживаемым. Кроме того, Kotlin полностью совместим с Java, что позволяет использовать существующие библиотеки и постепенно мигрировать на новый язык. Этот выбор позволяет создавать более эффективные, безопасные и современные приложения с минимальными усилиями.

В качестве среды разработки используется Android Studio. Она предоставляет все необходимые инструменты для создания, отладки и развертывания приложений для платформы Android. В ней можно создавать макеты пользовательского интерфейса, писать код на Kotlin, проводить отладку и тестирование приложения. \cite{refg}

Внутри Android Studio используется Gradle – это система управления зависимостями и сборкой проектов, используемая в Android Studio. С помощью Gradle будут определяться зависимости и настройки проекта.

Для удобства разработки используется система контроля версий Git, которая используется для отслеживания изменений в файловой системе проекта. Она позволяет сохранять историю изменений, возвращаться к предыдущим версиям проекта, совместно работать над кодом с другими участниками и управлять кодом в разных ветках разработки. \cite{refg}

Для написания интерфейса был выбран Jetpack Compose, так как это современный и мощный инструмент для создания пользовательских интерфейсов в Android. Compose позволяет описывать UI прямо в коде на Kotlin, что делает разработку более интуитивной и эффективной. Благодаря своей декларативной природе, Compose упрощает создание динамических и отзывчивых интерфейсов, а также поддерживает современные функции, такие как анимации и Material Design 3. Этот выбор позволяет создавать красивые и производительные приложения с минимальным количеством шаблонного кода. \cite{refh}

Для работы с базой данных был выбран Room, потому что это высокоуровневая библиотека, которая значительно упрощает взаимодействие с SQLite. \cite{reff} Room автоматизирует многие процессы, такие как создание таблиц, выполнение запросов и миграции, что позволяет сосредоточиться на логике приложения, а не на низкоуровневых деталях работы с базой данных. Благодаря интеграции с LiveData и Flow, Room также упрощает работу с данными в реальном времени, что делает его идеальным выбором для современных Android-приложений.

Для работы с сетевыми запросами был выбран Retrofit, так как это одна из самых популярных и удобных библиотек для взаимодействия с REST API. \cite{refg} Retrofit позволяет описывать API с помощью интерфейсов и аннотаций, что делает код чистым и легко читаемым. Благодаря своей интеграции с OkHttp, Retrofit обеспечивает высокую производительность и гибкость, а также поддерживает множество функций, таких как автоматическая сериализация/десериализация данных и обработка ошибок. Этот выбор позволяет эффективно работать с сетевыми запросами и сосредоточиться на бизнес-логике приложения.

Для тестирования сервера используется online-инструмент Postman, который позволяет отправлять HTTP-запросы к серверу и проверять ответы. \cite{refh} Он позволяет создавать и отправлять запросы различных типов (GET, POST, PUT, DELETE и т. д.), настраивать заголовки и параметры запроса, а также автоматизировать тестирование с помощью коллекций запросов и сценариев. Это будет полезно при работе с внешним API для интеграции с расписанием.
